\documentclass[11pt,a4paper]{article}
\usepackage{amsmath,amssymb}
\usepackage{graphicx}
\usepackage[numbers,sort]{natbib}
\usepackage{url}
\usepackage{color}
\usepackage{array}

% Page layout (manual, no geometry package)
\setlength{\textwidth}{16cm}
\setlength{\textheight}{24cm}
\setlength{\oddsidemargin}{0cm}
\setlength{\topmargin}{-1cm}
\setlength{\parindent}{0pt}
\setlength{\parskip}{6pt}

% Approximate booktabs
\newcommand{\toprule}{\hline\hline}
\newcommand{\midrule}{\hline}
\newcommand{\bottomrule}{\hline\hline}

% Custom macros
\newcommand{\method}{LoRA-Teleport}
\newcommand{\eg}{e.g.,~}
\newcommand{\ie}{i.e.,~}
\newcommand{\etal}{\textit{et al.}~}
\newcommand{\cossim}{\mathrm{cos\_sim}}
\newcommand{\PLACEHOLDER}[1]{\textcolor{red}{[PLACEHOLDER: #1]}}

\definecolor{red}{rgb}{0.8,0,0}

% =============================================================
\title{\textbf{Why Gradient Alignment Teleportation Fails in Sequential\\
       Continual Learning: A Mechanistic Study}}

\author{Anonymous Author(s) \\ \textit{Submitted for blind review}}

\date{}

\begin{document}
\maketitle
\thispagestyle{empty}

% =============================================================
\begin{abstract}
% =============================================================
We present a systematic empirical study of LoRA-based gradient alignment teleportation
for sequential continual learning (CL), adapting the recently proposed COST
framework~\citep{zhou2025cost} from multi-task to sequential settings.
Catastrophic forgetting remains a central challenge in CL, and gradient conflict between
old and new tasks is widely hypothesized as a proximal cause---yet direct manipulation
of gradient alignment during training has not been rigorously tested in the sequential regime.
We implement an online teleportation mechanism that detects gradient conflict at each
training step and applies LoRA perturbations to find loss-equivalent parameter configurations
with better gradient alignment.
Across six hyperparameter configurations on seq-CIFAR-10, all online teleportation variants
\emph{hurt} accuracy relative to a vanilla Experience Replay baseline ($-1.6\%$ to $-31.6\%$).
We identify a structural \emph{zero-order trap}: achieving loss invariance on one batch
requires large parameter perturbations ($\|\Delta\Theta\| = 25$--$30$) that catastrophically
disrupt the training trajectory---better loss invariance produces \emph{worse} accuracy.
A companion experiment further reveals that gradient conflict is \emph{ubiquitous} in 78.9\%
of training steps, providing insufficient discriminating signal to guide teleportation.
These findings constitute a rigorous negative result with a mechanistic explanation
of why LoRA teleportation fails in the sequential regime.
\end{abstract}

% =============================================================
\section{Introduction}
\label{sec:intro}
% =============================================================

Catastrophic forgetting---the rapid degradation of previously learned knowledge when a
neural network is trained on a new task---remains one of the core obstacles in continual
learning~\citep{kirkpatrick2017ewc, vandeven2022survey}.
A prominent hypothesis attributes forgetting to \emph{gradient conflict}: when the gradient
of the current task opposes that of old tasks, parameter updates that benefit new-task
performance necessarily degrade old-task performance~\citep{lopezpaz2017gem, yu2020pcgrad}.

This conflict hypothesis has motivated two main lines of work.
\emph{Constraint-based} methods (\eg GEM~\citep{lopezpaz2017gem}, A-GEM~\citep{chaudhry2019agem})
project current-task gradients to avoid increasing the replay-buffer loss.
\emph{Surgery-based} methods (\eg PCGrad~\citep{yu2020pcgrad}) modify gradient directions
to reduce inter-task conflict.
Both operate at the gradient level---altering how the optimizer steps.

A third possibility is \emph{parameter-space teleportation}: instead of modifying the gradient step,
find a loss-equivalent point in parameter space from which the gradient is naturally better aligned.
This idea was formalized by \citet{zhao2022teleport} and extended to multi-task learning (MTL) by
\citet{zhou2025cost} (COST), which uses low-rank adapters (LoRA) to efficiently find such points.
COST demonstrates strong results on simultaneous MTL benchmarks.

\textbf{This paper.}
We ask whether gradient alignment teleportation can reduce catastrophic forgetting in
\emph{sequential} CL, where old tasks are inaccessible and conflict is detected only
through a small replay buffer.
We implement an online variant (\method) and conduct a systematic sweep over
hyperparameters on seq-CIFAR-10.

\textbf{Contributions:}
\begin{itemize}
  \item \textbf{All six teleportation configurations hurt accuracy}
        ($-1.6\%$ to $-31.6\%$ vs.~ER baseline), ruling out hyperparameter explanations.
  \item \textbf{Zero-order trap:} the loss-invariance constraint forces
        $\|\Delta\Theta\| = 25$--$30$, disrupting the training trajectory through all
        non-target batches. Better loss invariance produces \emph{worse} accuracy.
  \item \textbf{Ubiquitous conflict:} 78.9\% of training steps exhibit
        $\cossim(g_\mathrm{old}, g_\mathrm{new}) < 0$.
        Pooled correlation with forgetting is $r = -0.25$ (weak, inconsistent across seeds),
        suggesting mini-batch $\cossim$ is not a reliable signal for teleportation decisions.
\end{itemize}

% =============================================================
\section{Background}
\label{sec:background}
% =============================================================

\paragraph{Sequential continual learning.}
In class-incremental learning (Class-IL), a model observes tasks $1, \ldots, T$ sequentially.
At test time it must classify across all seen classes with no task identity given.
Experience Replay (ER)~\citep{rolnick2019er} stores a buffer of past samples mixed with
current-task data. We use ER as the base method throughout.

\paragraph{Gradient conflict.}
For parameters $\theta$, conflict occurs when
$\cossim(g_i, g_j) < 0$ where $g_k = \nabla_\theta \mathcal{L}_k(\theta)$.
GEM~\citep{lopezpaz2017gem} enforces $g^\top g_k \geq 0$ for buffered tasks per step.
PCGrad~\citep{yu2020pcgrad} projects conflicting gradients onto each other's normal plane.

\paragraph{Symmetry teleportation.}
\citet{zhao2022teleport} exploit ReLU scaling symmetries to teleport to better-conditioned
points on the loss manifold while preserving the loss.
COST~\citep{zhou2025cost} extends this to LoRA perturbations in MTL,
minimizing gradient conflict while preserving per-task loss.

\paragraph{LoRA.}
Low-Rank Adaptation~\citep{hu2022lora} parameterizes $\Delta W = BA$ with
$B \in \mathbb{R}^{d \times r}$, $A \in \mathbb{R}^{r \times d}$, $r \ll d$.
After optimization, the LoRA is merged: $W \leftarrow W + BA$.

% =============================================================
\section{Method: Online LoRA Teleportation}
\label{sec:method}
% =============================================================

At each training step $s$ we draw a current-task batch $(x_\mathrm{new}, y_\mathrm{new})$
and an old-task batch $(x_\mathrm{old}, y_\mathrm{old})$ from a persistent replay iterator.
We compute gradient conflict via
\begin{equation}
  \cossim_s = \frac{g_\mathrm{old}^\top g_\mathrm{new}}{\|g_\mathrm{old}\|\|g_\mathrm{new}\| + \varepsilon}
  \label{eq:cossim}
\end{equation}
on the last $L=2$ layers (for efficiency). When $\cossim_s < \tau$, we optimize
rank-$r$ LoRA perturbations $\Delta\Theta = \sum_l B_l A_l$ via:
\begin{align}
  \mathcal{L}_\mathrm{teleport} &= \lambda_t \cdot \mathcal{L}_t
    - \cossim(\theta + \Delta\Theta)
    + \lambda_r \|\Delta\Theta\|_F^2
  \label{eq:teleport_obj} \\
  \mathcal{L}_t &= \big|\mathcal{L}(\theta{+}\Delta\Theta, x_\mathrm{old}) - \mathcal{L}(\theta, x_\mathrm{old})\big|
    + \big|\mathcal{L}(\theta{+}\Delta\Theta, x_\mathrm{new}) - \mathcal{L}(\theta, x_\mathrm{new})\big|
    \label{eq:lt}
\end{align}
where $\mathcal{L}_t$ is the \emph{loss invariance} term (zero means the perturbation is
loss-neutral on both batches) and $\lambda_t$ is its weight.
The cosine term in Eq.~\ref{eq:teleport_obj} requires second-order gradients.
After $n$ optimization steps, the LoRA is merged back into $\theta$.

% =============================================================
\section{Experiments: All Teleportation Variants Fail}
\label{sec:h7}
% =============================================================

\textbf{Setup.}
seq-CIFAR-10 (5 tasks, 2 classes/task), Class-IL protocol, ResNet-18.
Optimizer: SGD, lr=0.1, 10 epochs/task, seed=42.
ER baseline (buffer=500): 58.53\% Class-IL.
We sweep: check\_freq $\in \{1, 10\}$, $\lambda_t \in \{1, 10, 50\}$, $n \in \{5, 10, 20\}$;
rank-2 LoRA on the last 2 layers.

\begin{table}[t]
\centering
\caption{Online \method\ variants on seq-CIFAR-10 (Class-IL accuracy).
All variants underperform ER.
\textbf{E6 achieves the lowest avg\_lt yet the worst accuracy}---the zero-order trap.}
\label{tab:h7}
\begin{tabular}{|l|l|c|c|c|c|c|}
\toprule
Exp & Description & Class-IL $\uparrow$ & $\Delta$ & Triggers & avg\_lt $\downarrow$ & $\|\Delta\Theta\|$ \\
\midrule
E0 & ER (baseline)                     & \textbf{58.53\%} & ---        & 0    & ---  & ---   \\
\midrule
E5 & freq=10, $\lambda_t$=50, $n$=10   & 56.90\%          & $-1.63\%$  & 972  & 1.73 & 9--13 \\
E3 & freq=10, $\lambda_t$=1,  $n$=5    & 56.25\%          & $-2.28\%$  & 850  & 5.67 & 4--6  \\
E4 & freq=10, $\lambda_t$=10, $n$=10   & 53.51\%          & $-5.02\%$  & 927  & 1.92 & 9--13 \\
E2 & freq=1,  $\tau$=$-0.3$            & 48.87\%          & $-9.66\%$  & 1834 & 4.80 & ---   \\
E6 & freq=10, $\lambda_t$=10, $n$=20   & 48.22\%          & $-10.31\%$ & 1140 & \textbf{0.46} & 25--30 \\
E1 & freq=1,  $\lambda_t$=1,  $n$=5    & 26.93\%          & $-31.60\%$ & 4419 & ---  & ---   \\
\bottomrule
\end{tabular}
\end{table}

Table~\ref{tab:h7} shows all six configurations fall below ER.
The best variant (E5) is still $-1.63\%$ below baseline.
Increasing check\_freq to 1 is catastrophic ($-31.6\%$, 12$\times$ training slowdown).

\subsection{The Zero-Order Trap}

\textbf{E6 is the diagnostic case.}
With $\lambda_t = 10$ and $n=20$ steps, E6 achieves avg\_lt $= 0.46$---loss invariance
is working as intended.
Yet E6 produces the \emph{worst accuracy} among freq=10 variants ($-10.3\%$).

To satisfy $\mathcal{L}_t \approx 0$ on the target batch, the optimizer finds
$\|\Delta\Theta\| = 25$--$30$, compared to $4$--$6$ in the unconstrained regime.
This large perturbation is loss-neutral only at the specific target point $x^*$;
for any other batch $x \neq x^*$, the perturbation induces unpredictable loss changes.
With 3{,}130 batches per task, the training trajectory is corrupted by each teleportation event.

\begin{table}[h]
\centering
\caption{The zero-order trap: no regime simultaneously achieves small $\mathcal{L}_t$ and small $\|\Delta\Theta\|$.}
\label{tab:trap}
\begin{tabular}{|l|c|c|l|}
\toprule
Regime & avg\_lt & $\|\Delta\Theta\|$ & Net effect \\
\midrule
Unconstrained ($\lambda_t=1$)          & 5.67 & 4--6   & Target batch loss disrupted \\
Moderate ($\lambda_t=50$)              & 1.73 & 9--13  & Partial trajectory disruption \\
Strongly constrained ($\lambda_t=10$, $n$=20) & \textbf{0.46} & 25--30 & All other batches disrupted \\
\bottomrule
\end{tabular}
\end{table}

Table~\ref{tab:trap} shows there is no viable middle ground with rank-2 LoRA in 5--20 steps.
This is a structural constraint: the low-rank parameter space cannot simultaneously satisfy
the loss invariance constraint and find a small-norm $\Delta\Theta$ that improves alignment.

% =============================================================
\section{Analysis: Is Gradient Conflict the Right Target?}
\label{sec:h8}
% =============================================================

The failure of teleportation motivates a more fundamental question: is gradient conflict
a reliable signal for forgetting risk, or is it too ubiquitous and noisy to guide per-step decisions?

\textbf{Setup.}
We introduce a \emph{detect-only} mode: teleportation is disabled ($\tau=100$, never triggered),
while $\cossim_s$ is logged at every check step.
Three seeds (42, 123, 456) on seq-CIFAR-10 ER (10 epochs/task), check\_freq=10,
yielding 313 measurements per eligible task.
We compute per-task forgetting $f_t = \text{peak\_acc}_t - \text{final\_acc}_t$ and
Pearson $r$ between mean $\cossim_t$ per task and $f_t$.

\begin{table}[t]
\centering
\caption{Mean $\cossim(g_\mathrm{old}, g_\mathrm{new})$ during seq-CIFAR-10 ER training.
Parentheses: fraction of steps with $\cossim < 0$.}
\label{tab:h8_cossim}
\begin{tabular}{|l|c|c|c|c|c|}
\toprule
Seed & Task 1 & Task 2 & Task 3 & Task 4 & Mean neg. \\
\midrule
42  & $-0.173$ (70.6\%) & $-0.252$ (83.1\%) & $-0.214$ (77.0\%) & $-0.133$ (75.1\%) & 76.5\% \\
123 & $-0.179$ (73.2\%) & $-0.216$ (78.6\%) & $-0.304$ (89.1\%) & $-0.155$ (78.3\%) & 79.8\% \\
456 & $-0.234$ (80.8\%) & $-0.221$ (79.6\%) & $-0.268$ (84.7\%) & $-0.131$ (76.4\%) & 80.4\% \\
\midrule
\textbf{Mean} & $-0.195$ & $-0.230$ & $-0.262$ & $-0.140$ & \textbf{78.9\%} \\
\bottomrule
\end{tabular}
\end{table}

\textbf{Gradient conflict is ubiquitous.}
Table~\ref{tab:h8_cossim}: \textbf{78.9\% of all training steps exhibit gradient conflict}.
This is consistent with \citet{liu2026gradconflict}'s finding that 50--75\% of neurons
show conflicting gradient directions in LLM fine-tuning.
The mean $\cossim$ varies only from $-0.13$ to $-0.30$ across tasks and seeds
(a 0.17-unit range), leaving very little discriminating signal.

\textbf{Correlation is weak and inconsistent.}

\begin{table}[h]
\centering
\caption{Pearson $r$ between mean $\cossim$ per task and per-task forgetting.
Direct: $r(\bar{c}_t, f_t)$. Cross: $r(\bar{c}_t, f_{t-1})$.
Pooled uses 12 data points (3 seeds $\times$ 4 tasks).}
\label{tab:h8_corr}
\begin{tabular}{|l|c|c|l|}
\toprule
Seed & $r$ (direct) & $r$ (cross-task) & Verdict \\
\midrule
42  & $-0.43$ & $-0.73$ & Weakly supported \\
123 & $-0.01$ & $-0.32$ & Not supported    \\
456 & $-0.54$ & $+0.08$ & Mixed            \\
\midrule
\textbf{Pooled} & $\mathbf{-0.25}$ & --- & Inconclusive \\
\bottomrule
\end{tabular}
\end{table}

The pooled $r = -0.25$ (Table~\ref{tab:h8_corr}) is weakly negative but highly inconsistent.
Two confounds limit interpretation:
(1) \emph{Task-order confound}: later tasks always incur less forgetting and face slightly
    less conflict (fewer old tasks in buffer), creating a spurious negative correlation.
(2) \emph{Insufficient variance}: the 0.17-unit range of mean $\cossim$ provides too little
    signal to separate causal from confounded effects with only 4 tasks per seed.

% =============================================================
\section{Related Work}
\label{sec:related}
% =============================================================

\textbf{Replay-based CL.}
ER~\citep{rolnick2019er} and DER~\citep{buzzega2020der} mix past samples with current training.
We use ER as our base and augment it with teleportation.

\textbf{Gradient-based forgetting mitigation.}
GEM~\citep{lopezpaz2017gem} and A-GEM~\citep{chaudhry2019agem} provide \emph{first-order}
per-step guarantees by constraining gradient directions.
Our teleportation seeks a \emph{zero-order} batch-level guarantee.
The zero-order trap suggests first-order methods may be more robust: they modify only
the step direction, not the parameter value, and do not require compensating scale.

\textbf{Multi-task gradient surgery.}
PCGrad~\citep{yu2020pcgrad} projects conflicting gradients in the MTL setting.
MTL's simultaneous task access makes exact loss evaluation possible, avoiding the
replay-sampling noise that amplifies the zero-order trap.

\textbf{Symmetry teleportation.}
COST~\citep{zhou2025cost} evaluates loss on all tasks simultaneously in MTL;
our sequential setting must approximate old-task loss via a small replay buffer.
This approximation error interacts fatally with the zero-order constraint.

\textbf{Gradient similarity and forgetting.}
\citet{liu2026gradconflict} freeze conflicting neurons (50--75\% in LLMs) and reduce
forgetting by 59--82\%.
Their neuron-level approach is complementary to ours: freezing is a first-order-like
modification (don't update conflicting parameters), whereas teleportation is a zero-order
modification (change parameter values while preserving loss).
Our H8 result suggests that coarser mini-batch $\cossim$ does not have the resolution
needed for step-by-step teleportation decisions.

% =============================================================
\section{Limitations}
\label{sec:limitations}
% =============================================================

\textbf{Benchmark scope.} All experiments use seq-CIFAR-10 (5 tasks, ResNet-18).
Generalization to longer sequences, transformer architectures, or LLM fine-tuning is not confirmed.

\textbf{Low-rank budget.} Rank-2 LoRA with 5--20 steps may be insufficient.
Higher ranks might reduce $\|\Delta\Theta\|$---at substantially higher compute cost
(second-order gradients scale with rank).

\textbf{Statistical power.} With 4 eligible tasks per seed, correlation analysis has
12 pooled data points. Longer task sequences would give more reliable estimates.

\textbf{Alternative conflict signals.} Per-neuron conflict~\citep{liu2026gradconflict}
or accumulated trajectory cosine similarity might be more discriminative than per-step mini-batch $\cossim$.

% =============================================================
\section{Conclusion}
\label{sec:conclusion}
% =============================================================

We present a rigorous negative result: LoRA-based gradient alignment teleportation fails
in sequential continual learning. Two structural barriers account for the failure.

\textbf{The zero-order trap:} Loss invariance on a single batch forces
$\|\Delta\Theta\| = 25$--$30$, disrupting the training trajectory through all non-target batches.
The configuration with the best loss invariance (avg\_lt $= 0.46$) produces the worst accuracy.
This creates an inescapable dilemma: either enforce loss invariance and corrupt the trajectory,
or relax it and allow the optimizer to sacrifice it for $\cossim$ gain.

\textbf{Ubiquitous conflict:} 78.9\% of training steps show gradient conflict.
The low variance in mean $\cossim$ per task (0.17 units, pooled $r = -0.25$, high seed variance)
shows that mini-batch $\cossim$ cannot reliably guide per-step teleportation.

These findings do not contradict the gradient conflict hypothesis at the neuron level.
They identify \emph{parameter-space teleportation} as an ineffective vehicle for conflict
reduction in the sequential CL regime.
Future work might explore first-order gradient projection for sequential CL,
or neuron-level selective approaches inspired by~\citet{liu2026gradconflict},
as alternatives that avoid the zero-order trap.

\bibliographystyle{plainnat}
\bibliography{bibliography}

\end{document}
